% Chapter 9 worked example. Included by ch09_resources.tex.
\section{Worked Example: Vendor Procurement and Renegotiation}
\label{sec:comp-ch09-worked-example}

\subsection{Purpose and Status}
This worked example demonstrates how the methods in Chapter~9 can be
translated into a reviewable institutional decision. All numerical inputs are
synthetic unless a source is identified explicitly. They are intended to show the
logic of the analysis, not to report empirical findings or establish a universal
threshold. Local use requires replacement of the assumptions with current,
versioned evidence and approval through the relevant clinical, managerial,
economic, legal, technical, and governance processes.

A hospital is procuring a clinically consequential AI platform. Three vendors differ in evidence quality, interoperability, pricing, monitoring access, update control, and exit support. A prior procurement also requires renegotiation because implementation costs and vendor restrictions exceeded the original assumptions. 

\subsection{Decision Problem and Comparators}
\textbf{Decision problem.} Which vendor, if any, should proceed to a controlled contract, and what remedies or exit actions are required for the existing platform?

The comparator set is deliberately broader than the existing pathway. This prevents
a new technology from receiving a lower evidentiary threshold than staffing,
workflow, shared-service, conventional digital, or no-change alternatives.

\begin{longtable}{@{}p{0.08\textwidth}p{0.84\textwidth}@{}}
\caption{Comparators for the Chapter 9 worked example}
\label{tab:comp-ch09-comparators}\\
\toprule
\textbf{Option} & \textbf{Comparator or strategy} \\
\midrule
\endfirsthead
\toprule
\textbf{Option} & \textbf{Comparator or strategy} \\
\midrule
\endhead
\bottomrule
\endlastfoot
1 & Vendor A: strong evidence, high fixed price \\
2 & Vendor B: lower price, limited monitoring access \\
3 & Vendor C: shared-service model with stronger exit terms \\
4 & Internal or non-AI process alternative \\
5 & Renegotiate, replace, restrict, or retire the existing platform \\
\end{longtable}

\subsection{Model and Decision Structure}
The analytical structure should follow the decision rather than the available
software. The team should map the causal pathway from intervention input to
information, human decision, action, outcome, resource consequence, and review.
The structure should also identify capacity constraints, uncertainty, subgroup
variation, implementation requirements, and events that would change the
recommendation.

A compact quantitative relationship used in this example is:
\begin{equation}
S_v=\sum_{d=1}^{D}w_ds_{vd},\qquad G_v=\prod_{k=1}^{K}g_{vk}
\label{eq:comp-ch09-key-relationship}
\end{equation}
The equation is an organizing aid rather than a substitute for disaggregated evidence,
qualitative findings, or mandatory safeguards. Its variables and interpretation must
be defined in the local protocol.

\subsection{Synthetic Parameters and Evidence Sources}
Table~\ref{tab:comp-ch09-parameters} provides the base-case inputs. A
production analysis should add parameter distributions, correlations, structural
alternatives, data dates, system versions, and evidence-confidence ratings.

\begin{longtable}{@{}p{0.32\textwidth}p{0.22\textwidth}p{0.38\textwidth}@{}}
\caption{Synthetic parameters for the Chapter 9 worked example}
\label{tab:comp-ch09-parameters}\\
\toprule
\textbf{Parameter} & \textbf{Base value} & \textbf{Source or interpretation} \\
\midrule
\endfirsthead
\toprule
\textbf{Parameter} & \textbf{Base value} & \textbf{Source or interpretation} \\
\midrule
\endhead
\bottomrule
\endlastfoot
Mandatory safeguard domains & 7 & Procurement specification \\
Weighted domains & 7 & Local scorecard \\
Expected annual volume & 180,000 uses & Activity forecast \\
Maximum acceptable variable exposure & EUR 420,000/year & Budget condition \\
Required update notice & 90 days & Contract requirement \\
Minimum data-export period after termination & 12 months & Continuity requirement \\
Remediation period for existing vendor & 120 days & Renegotiation proposal \\
\end{longtable}

\subsection{Step-by-Step Analysis}
The sequence below is intended to make the analysis reproducible and to ensure
that evidence, economics, governance, and implementation develop together.

\begin{longtable}{@{}p{0.07\textwidth}p{0.55\textwidth}p{0.30\textwidth}@{}}
\caption{Analytical sequence}
\label{tab:comp-ch09-analysis-steps}\\
\toprule
\textbf{Step} & \textbf{Required work} & \textbf{Decision record} \\
\midrule
\endfirsthead
\toprule
\textbf{Step} & \textbf{Required work} & \textbf{Decision record} \\
\midrule
\endhead
\bottomrule
\endlastfoot
1 & Define the service problem, comparator set, intended and excluded uses, and measurable outcomes. & Evidence and responsible owner to be recorded \\
2 & Require version-matched evidence and complete technical, data, security, economic, and workflow due diligence. & Evidence and responsible owner to be recorded \\
3 & Apply mandatory safeguards before weighted comparison. & Evidence and responsible owner to be recorded \\
4 & Model pricing and lifecycle cost under realistic volume, update, support, and exit scenarios. & Evidence and responsible owner to be recorded \\
5 & Convert monitoring, audit, update, continuity, portability, and exit requirements into enforceable clauses. & Evidence and responsible owner to be recorded \\
6 & For the existing platform, compare prospective continuation, renegotiation, restriction, replacement, and decommissioning without allowing sunk cost to determine the decision. & Evidence and responsible owner to be recorded \\
\end{longtable}

\subsection{Illustrative Outputs}
The outputs in Table~\ref{tab:comp-ch09-outputs} show how technical,
clinical, operational, economic, equity, and governance findings should be kept
visible rather than compressed into a single favourable or unfavourable score.

\begin{longtable}{@{}p{0.30\textwidth}p{0.24\textwidth}p{0.38\textwidth}@{}}
\caption{Illustrative model and decision outputs}
\label{tab:comp-ch09-outputs}\\
\toprule
\textbf{Output} & \textbf{Result} & \textbf{Interpretation} \\
\midrule
\endfirsthead
\toprule
\textbf{Output} & \textbf{Result} & \textbf{Interpretation} \\
\midrule
\endhead
\bottomrule
\endlastfoot
Vendor B weighted score & Highest before safeguards & Ineligible because monitoring access fails a mandatory requirement \\
Vendor A evidence profile & Strongest & Price and exit terms require negotiation \\
Vendor C continuity profile & Strongest & Local evidence and governance arrangements require confirmation \\
Existing vendor decision & Time-limited remediation with parallel exit plan & Continuation is conditional \\
Preferred award approach & Competitive dialogue followed by conditional contract & No automatic award from aggregate score \\
\end{longtable}

\subsection{Sensitivity, Uncertainty, and Subgroups}
At minimum, the completed analysis should vary the parameters most likely to
change the decision, test at least one credible alternative model structure, and
report subgroup results separately where performance, access, response, burden,
or opportunity cost may differ. Deterministic analysis should identify thresholds;
probabilistic analysis should be used where credible parameter distributions and
correlations are available. Scenario analysis is preferable when structural or
future uncertainty cannot be represented defensibly by one probability model.

The record should distinguish uncertainty that can be reduced through evidence,
uncertainty that can be managed through safeguards, and uncertainty that remains
too consequential to accept. Missing evidence should not be represented as an
average or neutral value without justification.

\subsection{Interpretation and Recommendation}
Proceed only with a vendor that passes every mandatory safeguard and accepts version-specific evidence, monitoring access, transparent pricing, material-change control, continuity, and exit obligations. Continue the existing platform only under a time-limited remediation agreement and maintain a funded replacement plan.

The recommendation is conditional rather than permanent. It applies only to the
specified population, setting, intervention configuration, evidence cut-off, and
version. The following conditions should be incorporated into the approval,
contract, implementation, monitoring, or renewal record:

\begin{itemize}
 \item mandatory safeguards are evaluated independently of weighted preferences.
 \item missing evidence receives clarification, conservative treatment, or exclusion.
 \item contracted pricing includes integration, monitoring, updates, incident support, and exit.
 \item the institution retains audit, data, log, and publication rights needed for accountability.
 \item renewal is linked to realized value and current evidence rather than sunk cost.
\end{itemize}

\subsection{Decision Statement}
\begin{quote}
For the specified decision problem and comparator set, the institution should
\emph{[proceed / proceed with conditions / redesign / pilot / restrict / defer /
replace / retire]}. The decision applies to \emph{[population, setting, version,
and evidence period]}, is owned by \emph{[accountable authority]}, and will be
reviewed on \emph{[date]} or earlier if \emph{[trigger events]} occur. The
evidence, assumptions, unresolved uncertainty, mandatory safeguards, monitoring
thresholds, and exit arrangements are recorded in the accompanying dossier.
\end{quote}
